\begin{proposition}[Root-atlas base change and relative duality]\hypertarget{tag-033I}{}\label{push:src:E.3.1}
\StatementRecord{push:src:E.3.1}{proved}{033I}{the normalized torus root charts and the periodic triangular cusp charts; \cite[Chapter III, Theorem 4.12(ii), pp.~134--135, and Theorem 4.10, pp.~132--133]{BS76}; \cite[introduction, equation (1), p.~261; Section 1, pp.~264--265; Proposition 2.1, p.~266]{RRV71}}

For every root lift $\bar f:X\to\mathfrak B_{k,l}=\mathbb P(3k,4l)$
of the retained coprime-multiplicity construction, the direct images
$R^q\bar f_*\mathcal O_X$, for $q=0,1,2$, are locally free of
ranks $1,2,1$. The canonical geometric-point maps
\[
 (R^q\bar f_*\mathcal O_X)\otimes k(b)
 \longrightarrow H^q(X_b,\mathcal O_{X_b})
\]
are isomorphisms, and the trace induces the canonical isomorphism
\[
 (R^2\bar f_*\mathcal O_X)^\vee
 \simeq\bar f_*\omega_{X/\mathfrak B_{k,l}}.
\]
Both comparisons commute with the analytic root-atlas groupoid.
Here the geometric fibre at a finite root point is the smooth torus
on the normalized root chart.
\end{proposition}

\begin{proof}
The fibre-divisor identities $f^*p_1=3kS_1$ and $f^*p_2=4lS_2$
define the root lift. On a full root chart its atlas pullback is
\[
 [J'_j/\mu_{m_j}]
 \mathbin{\times_{[\Delta_{s_j}/\mu_{m_j}]}}
 \Delta_{s_j}=J'_j,
 \qquad (m_1,m_2)=(3k,4l).
\]
The affine quotient action is free and the family $J'_j\to\Delta_{s_j}$
is the original smooth proper torus family. The atlas map is etale.
At the cusp the unchanged local charts are $q=z_0z_1z_2$, with the
two-factor and smooth charts on the other strata. They have smooth
total space and pure two-dimensional Cartier fibres. A base parameter
is a non-zero-divisor, so the atlas family is flat. Properness is
properness of the compact torus and cusp fillings. These local
comparisons make the atlas families proper flat Gorenstein morphisms
of relative dimension two.

A smooth two-torus has structure-sheaf dimensions $(1,2,1)$.
For the cusp fibre $W$, its normalization resolution has terms
$\mathcal O_{\mathrm dP_6}$, three copies of
$\mathcal O_{\mathbb P^1}$, and two point sheaves. Its global
section complex, in the fixed branch orientations, is
\[
 \mathbb C\xrightarrow{0}\mathbb C^3
 \xrightarrow{(a,b,c)\mapsto(a-b+c,a-b+c)}\mathbb C^2.
\]
The resolution is exact by gluing functions on the coordinate
branches of the product charts. The positive-degree cohomology of
its rational strata vanishes. Hence its cohomology has dimensions
$(1,2,1)$ as well. Its cup product is the two-triangle torus product;
the cocycles $(0,-1,-1)$ and $(1,0,-1)$ pair to a generator in
degree two. In particular the alternating cup pairing is
nondegenerate. Thus the three dimensions are constant on every
atlas chart, including its centre.

The constant-cohomology theorem
\cite[Chapter III, Theorem 4.12(ii), pp.~134--135]{BS76}
now gives local freeness, and the canonical analytic base-change
maps are those of
\cite[Chapter III, Theorem 4.10, pp.~132--133]{BS76}.
The hypotheses are properness, flatness of the structure sheaf,
a reduced smooth base, and the constant fibre dimensions just
computed. The original constant-dimension formulation is
\cite[Section 7, Satz 5, pp.~63--64]{Grauert1960D}, with
\cite[pp.~35--36]{Grauert1963ErratumD}.
On a double atlas overlap, both pullbacks identify with the direct
image of the same pulled-back family. Functoriality makes these
identifications satisfy the triple-overlap cocycle condition, so
they descend with the original stabilizer representations.

For an atlas family $p:Y\to U$, proper analytic duality gives
\[
 Rp_*\omega_{Y/U}[2]
 \simeq R\mathcal Hom_U(Rp_*\mathcal O_Y,\mathcal O_U).
\]
This is the canonical trace identity of
\cite[introduction, equation (1), p.~261; Section 1, pp.~264--265;
Proposition 2.1, p.~266]{RRV71}.
The smooth-space dualizing complex is $K_Y[3]$ and the relative
Gorenstein complex is $\omega_{Y/U}[2]$; see
\cite[Section 3, Theorem 1, p.~81; Proposition 2 and the
complete-intersection calculation, p.~91]{RR70}.
The locally free cohomology sheaves already obtained have no
positive sheaf Ext against $\mathcal O_U$. Degree $-2$ of the
trace identity is therefore
$p_*\omega_{Y/U}\simeq(R^2p_*\mathcal O_Y)^\vee$.
The trace commutes with the local biholomorphisms on the atlas
and has the same cocycle compatibility. It consequently descends
to the asserted isomorphism on $\mathfrak B_{k,l}$.
\end{proof}

